\chapter{Metodologia}

Este capítulo apresenta o desenho de pesquisa adotado para investigar como bares e restaurantes de Porto Alegre que adotaram cardápios digitais com autoatendimento desde 2020 percebem os benefícios, riscos e desafios dessa transição, e como seus clientes vivenciam essa interação. Em linha com a natureza do fenômeno --- que envolve tanto decisões organizacionais quanto experiências individuais ---, opta-se por uma abordagem de \textbf{métodos mistos}, combinando um estudo de caso múltiplo de natureza qualitativa com um \emph{survey} quantitativo aplicado a frequentadores. As próximas seções caracterizam a pesquisa, descrevem suas etapas, detalham os instrumentos de coleta, apresentam os critérios de seleção dos casos, esboçam o plano de análise e tratam dos aspectos éticos.

\section{Caracterização da pesquisa}
\label{sec:caracterizacao}

Quanto à \textbf{natureza}, esta pesquisa é classificada como aplicada, pois visa gerar conhecimento de uso prático para gestores de bares e restaurantes, projetistas de soluções de cardápio digital e formuladores de políticas setoriais. Quanto aos \textbf{objetivos}, é descritiva-exploratória: descritiva porque busca caracterizar como a transição para cardápios digitais ocorre nos estabelecimentos investigados; exploratória porque, apesar do crescimento do fenômeno no pós-pandemia, ainda há poucos estudos brasileiros que articulam, no mesmo recorte, a perspectiva dos gestores e a dos clientes em estabelecimentos não pertencentes a grandes redes.

Quanto à \textbf{abordagem}, adota-se uma estratégia de métodos mistos, com prioridade sobre a fase qualitativa e papel complementar para a fase quantitativa. Seguem-se as diretrizes de \citeonline{venkatesh2013bridging} para condução e validação de pesquisas mistas em Sistemas de Informação, em especial: (i) a explicitação dos motivos para a mistura --- aqui, triangular as percepções declaradas pelos gestores com a experiência relatada por clientes ---; (ii) a definição do \emph{design} --- neste trabalho, \emph{concorrente convergente paralelo}, com coletas qualitativa e quantitativa conduzidas em paralelo, peso analítico maior atribuído à fase qualitativa e integração das duas fontes no momento da interpretação dos resultados; e (iii) a avaliação separada da qualidade de cada componente (qualitativa, quantitativa e meta-inferencial).

A escolha por métodos mistos é também consistente com trabalhos brasileiros recentes sobre adoção de tecnologia no setor de alimentação fora do lar. \citeonline{fioreti2026aplicativos} combinam estudo de caso com \emph{survey} em estabelecimentos do setor; \citeonline{costa2021implicacoes} valem-se de entrevistas semiestruturadas para captar implicações da transformação digital em pequenos negócios alimentícios. Ambos os trabalhos demonstram a viabilidade e a riqueza analítica de articular dados qualitativos e quantitativos no recorte deste TCC.

Quanto aos \textbf{procedimentos técnicos}, esta pesquisa é organizada como estudo de caso múltiplo holístico, conforme \citeonline{yin2015estudo}, em que cada estabelecimento selecionado constitui uma unidade de análise. O estudo de caso múltiplo permite contrastar configurações distintas e fortalecer a validade externa via replicação literal e teórica \cite{yin2015estudo}. O \emph{survey} com frequentadores, por sua vez, fornece um pano de fundo sobre a experiência típica do consumidor com cardápios digitais por \emph{QR Code} em Porto Alegre, contra o qual os achados dos casos podem ser interpretados. A classificação quanto à natureza, aos objetivos, à abordagem e aos procedimentos segue os critérios apresentados por \citeonline{gil2017projetos}.

\section{Etapas da pesquisa}
\label{sec:etapas}

A pesquisa é estruturada em duas grandes fases, alinhadas à divisão entre TCC I e TCC II do curso. A \textbf{Fase 1} (TCC I) compreende: revisão de literatura nos eixos temáticos delimitados na fundamentação (transformação digital no \emph{food service}; cardápios digitais e \emph{QR Code}; autoatendimento e \emph{UX} em \emph{mobile}; adoção tecnológica via TAM e UTAUT/UTAUT2; LGPD aplicada a restaurantes); definição do problema e dos objetivos; elaboração e refinamento dos instrumentos (roteiro de entrevista no Apêndice~\ref{ap:roteiro} e questionário no Apêndice~\ref{ap:questionario}); e produção dos termos éticos (Apêndice~\ref{ap:termos}).

A \textbf{Fase 2} (TCC II) compreende: prospecção e contato com estabelecimentos candidatos; agendamento e condução das entrevistas com gestores; lançamento e divulgação do \emph{survey} com frequentadores; transcrição, codificação e análise das entrevistas; tabulação e análise descritiva das respostas do \emph{survey}; triangulação entre as duas fontes; redação dos capítulos de resultados e discussão; e defesa. O cronograma detalhado dessa fase está no Capítulo~\ref{ch:cronograma}.

\section{Instrumentos de coleta}
\label{sec:instrumentos}

\subsection{Entrevista semiestruturada com gestores}
\label{sec:entrevista}

A coleta qualitativa será conduzida por meio de \textbf{entrevistas semiestruturadas} com gestores ou proprietários responsáveis pela decisão sobre o cardápio digital em cada estabelecimento selecionado. A entrevista é planejada para uma duração entre 45 e 60 minutos, será gravada em áudio mediante consentimento expresso por meio do Termo de Consentimento Livre e Esclarecido (Apêndice~\ref{ap:termos}) e poderá ser realizada presencialmente ou por videoconferência, conforme a disponibilidade do participante.

O roteiro completo está apresentado no Apêndice~\ref{ap:roteiro} e organiza-se em seis blocos modulares:

\begin{itemize}
    \item \textbf{Bloco 0 --- Abertura}: contextualiza a pesquisa e formaliza o consentimento;
    \item \textbf{Bloco 1 --- Caracterização do estabelecimento e do entrevistado}: porte, anos de operação, perfil do tomador de decisão;
    \item \textbf{Bloco 2 --- Trajetória de adoção}: ano e gatilho da adoção do cardápio digital, decisões anteriores e posteriores à pandemia;
    \item \textbf{Bloco 3 --- Implementação e operação}: tecnologia adotada, fornecedor, processos internos, integrações;
    \item \textbf{Bloco 4 --- Percepção sobre clientes}: aceitação, reclamações recorrentes, pedidos de cardápio físico;
    \item \textbf{Bloco 5 --- Benefícios, riscos e desafios}: ganhos operacionais percebidos, riscos jurídicos e de dados, perspectiva futura.
\end{itemize}

A estrutura modular permite adaptar a profundidade de cada bloco ao perfil do estabelecimento (descrito na Seção~\ref{sec:criterios}) e ao tempo disponível do entrevistado. Esse formato segue a tradição de trabalhos brasileiros sobre transformação digital em pequenos negócios alimentícios, como \citeonline{costa2021implicacoes}, e dialoga com o estudo exploratório de \citeonline{martiniano2022delivery}, que conduziu 12 entrevistas em profundidade para mapear motivações e percepções relacionadas a aplicativos de \emph{delivery}.

\subsection{Questionário com frequentadores}
\label{sec:questionario}

A coleta quantitativa será realizada por meio de um \textbf{questionário autopreenchível}, hospedado em formulário \emph{web} (\emph{Google Forms}), com duração estimada de 5 a 8~minutos. A estrutura completa está no Apêndice~\ref{ap:questionario} e contempla quatro seções:

\begin{enumerate}
    \item \textbf{Perfil sociodemográfico} (idade, gênero, escolaridade, bairro de residência ou frequência);
    \item \textbf{Hábitos de frequência} em bares e restaurantes (frequência, tipos de estabelecimento, gasto médio);
    \item \textbf{Experiência com cardápio digital via \emph{QR Code}} (utilidade percebida, facilidade de uso, motivação hedônica, comparação com cardápio físico, situações de fricção);
    \item \textbf{Atitudes e intenções} (preferência geral, intenção de uso futuro, opiniões livres em campo aberto).
\end{enumerate}

A construção das questões da Seção~3 do questionário inspira-se nas construções da UTAUT2 \cite{venkatesh2012consumer}, especialmente \emph{expectativa de performance}, \emph{expectativa de esforço}, \emph{motivação hedônica}, \emph{valor (custo-benefício)} e \emph{hábito}, por sua adequação ao contexto de uso voluntário pelo consumidor final. A meta é obter pelo menos 80 respostas válidas. A divulgação ocorrerá por (i) redes sociais e grupos de interesse de Porto Alegre, (ii) eventual \emph{QR Code} disponibilizado nos estabelecimentos participantes, mediante autorização do gestor, e (iii) divulgação institucional na rede da UNISINOS. Recorte temporal e geográfico do respondente são checados como filtros iniciais para garantir aderência ao escopo. O método quantitativo via \emph{Google Forms} é comparável ao adotado por \citeonline{costa2022comportamento} no estudo da percepção de consumidores brasilienses sobre cardápios digitais.

\section{Critérios de seleção dos estabelecimentos}
\label{sec:criterios}

A amostragem dos estabelecimentos é \textbf{intencional e teoricamente orientada}, com o objetivo de contrastar configurações que evidenciem o papel da pandemia como inflexão na adoção tecnológica. Pretende-se selecionar \textbf{três estabelecimentos}, distribuídos em até três bairros de Porto Alegre (preferencialmente Cidade Baixa, Moinhos de Vento e Centro), atendendo aos três perfis a seguir:

\begin{description}
    \item[Perfil A.] \textbf{Pré-2020 que adotou.} Estabelecimento em operação antes de março de 2020 que, durante ou após a pandemia, adotou cardápio digital com autoatendimento.
    \item[Perfil B.] \textbf{Pré-2020 que resistiu.} Estabelecimento em operação antes de março de 2020 que, deliberadamente, manteve o cardápio físico ou apenas digital sem autoatendimento.
    \item[Perfil C.] \textbf{Pós-2020 nascido digital.} Estabelecimento aberto a partir de 2020 que já nasceu operando com cardápio digital com autoatendimento.
\end{description}

A composição dos perfis dialoga com o argumento de \citeonline{giovanini2024pandemia} de que a difusão de tecnologias digitais no \emph{food service} desencadeada pela pandemia tem efeitos persistentes e estruturais, o que justifica olhar tanto para quem mudou de prática quanto para quem optou explicitamente por não mudar e para quem nasceu já com a nova configuração. Como critérios de \textbf{inclusão}, adotam-se: porte pequeno ou médio, com tomador de decisão acessível; estabelecimento independente ou rede local pequena. Como critérios de \textbf{exclusão}: redes nacionais ou internacionais franqueadas (cujas decisões tecnológicas são centralizadas), restaurantes de \emph{fine dining} de altíssimo \emph{ticket} médio e operações exclusivamente de \emph{delivery} (\emph{dark kitchens}), por configurarem realidades operacionais distintas do recorte.

\section{Plano de análise}
\label{sec:analise}

O material qualitativo será tratado por meio de \textbf{análise de conteúdo} de orientação categorial, conforme \citeonline{bardin2016analise}, com as etapas clássicas de pré-análise, exploração do material, tratamento dos resultados e interpretação. As entrevistas serão transcritas integralmente, e as categorias iniciais derivam dos blocos do roteiro, sendo refinadas indutivamente a partir do conteúdo emergente das transcrições. Para apoiar a codificação, será utilizada planilha estruturada ou \emph{software} de análise qualitativa (a definir).

O material quantitativo do \emph{survey} será analisado por \textbf{estatística descritiva}, com cálculo de frequências, médias e medianas por questão, e por cruzamentos relevantes (por exemplo, faixa etária × intenção de uso, frequência de ida × percepção de utilidade). A correlação entre construções da UTAUT2, quando aplicável, é prevista como análise complementar; em função do tamanho amostral previsto (80+ respostas), as análises serão tratadas como exploratórias, não inferenciais.

A \textbf{integração das duas fontes} (triangulação) seguirá o modelo de meta-inferências discutido por \citeonline{venkatesh2013bridging}: para cada categoria emergente das entrevistas, verifica-se se há evidência convergente, divergente ou complementar nos dados do \emph{survey}. Achados convergentes reforçam a explicação; achados divergentes são tratados como pontos de aprofundamento na discussão.

\section{Aspectos éticos}
\label{sec:etica}

Todos os participantes da pesquisa --- gestores e frequentadores --- assinam ou anuem ao Termo de Consentimento Livre e Esclarecido (TCLE), cujo modelo está no Apêndice~\ref{ap:termos}. O TCLE descreve os objetivos da pesquisa, os procedimentos de coleta, os riscos previstos (mínimos), a forma de armazenamento dos dados, o direito de desistência a qualquer momento e os contatos do pesquisador e do orientador.

Para os gestores, aplica-se ainda o \textbf{Termo de Confidencialidade} (modelo institucional UNISINOS, anexado ao repositório do projeto), garantindo que informações comerciais sensíveis discutidas na entrevista possam ser anonimizadas no relatório final, quando assim solicitado. Quando houver registro de imagens ou da fachada dos estabelecimentos, será aplicado o \textbf{Termo de Autorização de Uso de Imagem} também previsto pela instituição.

O tratamento de dados pessoais coletados pelo \emph{survey} respeita as exigências da Lei Geral de Proteção de Dados (LGPD --- Lei~13.709/2018) \cite{martins2021principiologia,almeida2022impactos}: coleta limitada à finalidade declarada, armazenamento restrito a sistemas controlados pelo pesquisador, ausência de coleta de dados sensíveis e descarte após o prazo necessário à pesquisa. A discussão jurídica específica sobre cardápios digitais e relações de consumo dialoga com \citeonline{boscolo2021lgpd} e \citeonline{rocha2025abusividade}.

O projeto envolve pesquisa com seres humanos cujo objeto não compreende intervenção biomédica, coleta de dados sensíveis ou populações vulneráveis, e cujos riscos previstos são de natureza social mínima. A necessidade de submissão a Comitê de Ética em Pesquisa (CEP) será verificada junto às diretrizes institucionais vigentes da UNISINOS, em conjunto com o orientador, antes do início da coleta no TCC II; caso a submissão se mostre necessária ou recomendada, será providenciada com a antecedência adequada em relação ao cronograma de coleta.
